\documentclass[12pt]{article}
\usepackage[T2A]{fontenc}
\usepackage[utf8]{inputenc}
\usepackage[russian]{babel}
\usepackage{geometry}
\geometry{margin=2cm}
\begin{document}
\section*{Условия задач из презентации \texttt{Original\_stereometry\_tasks.pptx}}
Текст перепроверен: извлечение учитывает формульные элементы (в том числе корни, индексы и дроби).
\begin{enumerate}
  \item Точка G лежит на боковом ребре SC правильной шестиугольной пирамиды SABCDEF свершиной S.
  \par а) Постройте точку пересечения прямой BG с плоскостью боковой грани ESF.
  \par б) Найдите угол между прямыми BG и AD,если стороны основания пирамиды равны 6,боковые ребра равны 3√(13),а SG:GC=1:2.
  \item Основание пирамиды SABCD – параллелограмм ABCD.Точка K – середина ребра SD.
  \par а) Плоскость проходит через точку K параллельно медианам BM и SN граней BSC и ASD.Постройте прямую пересечения этой плоскости с плоскостью основания пирамиды.
  \par б) Найдите угол между прямыми BM и SN,если пирамида SABCD – правильная, причем все ее ребра равны.
  \item Дана прямая призма ABCA\_1B\_1C\_1.Плоскость, проходящая через центр основания A\_1B\_1C\_1 и середину K ребра BC,параллельна прямой AB.Эта плоскость пересекает прямую CC\_1 в точке L.
  \par а) Докажите, что CL=3CC\_1.
  \par б) Найдите угол между прямыми KL и AC\_1,если∠ACB=90°и AA\_1=AC=(1)/(4)BC.
  \item Основание пирамиды ABCD−прямоугольныйтреугольник ABC. Высота пирамиды проходит через середину гипотенузы AB.
  \par а) Докажите, чтобоковые ребра пирамиды образуют равные углы с плоскостью основания.
  \par б) Известно, что BC:AC=√(3):1,а угол между боковой гранью BDC и плоскостью основания равен 60°.Найдите углы двух других боковых граней с плоскостью основания.
  \item Вправильной треугольной призме ABCA\_1B\_1C\_1 стороны основания равны 5, боковые ребра равны 2, точка D−середина ребра CC\_1.
  \par а) Постройте прямую пересечения плоскостей ABC и ADB\_1.
  \par б) Найдите угол между плоскостями ABC и ADB\_1.
  \item Воснованиипрямой призмы ABCDA\_1B\_1C\_1D\_1 лежит квадрат ABCD со стороной 2, а высота призмы равна 1. Точка E лежит на диагонали BD\_1,причем BE=1.
  \par а) Постройтесечение призмы плоскостью A\_1C\_1E.
  \par б) Найдите угол между плоскостью сечения и плоскостью ABC.
  \item Основание пирамиды SABCD−равнобедренная трапеция ABCD с основаниями AD и BC,причем AD=2BC=2AB.Высота SH пирамиды проходит через точку пересечения прямых AB и CD.
  \par а) Докажите, что треугольник SBD прямоугольный.
  \par б) Найдите расстояние от точки C до плоскости ASD,если SH=BC=4.
  \item Боковые ребрапирамиды SABC с вершиной S попарно перпендикулярны.
  \par а) Докажите, что высота SH пирамиды проходит через точку пересечения высот основания ABC.
  \par б) Найдите SH,если боковые ребра равны 2,2,и 7√(2).
  \item Основание шестиугольной пирамиды SABCDEF−правильный шестиугольник ABCDEF. Высота пирамиды втрое больше стороны основания и проходит через точку E.
  \par а) Докажите, чтоугол между боковой гранью ASB и плоскостью основания равен 60°.
  \par б) Найдите расстояние от точки C до плоскости ASB,если сторона основания пирамиды равна 4.
  \item Точка K лежит на стороне AB основания ABCD правильной четырехугольной пирамиды SABCD,все ребра которой равны. Плоскостьαпроходит через точку K параллельно плоскости ASD.Сечение пирамиды плоскостьюα- четырехугольник, в который можно вписать окружность.
  \par а) Докажите, что BK=2AK.
  \par б) Найдите расстояние от вершины S до плоскостиα, если все ребра пирамиды равны 1.
  \item Данатреугольная пирамида SABC;O−точка пересечения медиан основания ABC.
  \par а) Докажите, что плоскость, проходящая через прямую AB и середину M ребра SC,делит отрезок SO в отношении 3:1, считая от вершины S.
  \par б) Найдитеугол между прямой BC и плоскостью ABM,если пирамида правильная, а угол между прямой, проходящей через точку M и середину ребра AB,и прямой SO равен 45°.
  \item Данаправильная четырехугольная пирамида SABCD с вершиной S. Все ребра пирамиды равны. Точка M-середина ребра BC.
  \par а) Докажите, что ортогональная проекция середины ребра AB на плоскость CSD делит медиану SN этой грани в отношении 1:2,считая от вершины S.
  \par б) Найдитеугол между прямой SM и плоскостью CSD.
  \item Основания ABC и A\_1B\_1C\_1 призмы ABCA\_1B\_1C\_1- равносторонние треугольники. Отрезок, соединяющий центр O основания ABC призмы с вершиной C\_1,перпендикулярен основаниям призмы.
  \par а) Докажите, что плоскости ABC\_1 и OCC\_1 перпендикулярны.
  \par б) Найдитеугол между прямой AA\_1 и плоскостью ABC\_1,если боковое ребро призмы равно стороне основания.
  \item Основаниепирамиды SABCD−ромб ABCD с углом 60°при вершине A.Боковое ребро SD перпендикулярно плоскости основания и равно стороне основания.
  \par а) Докажите, что прямые AC и SB перпендикулярны.
  \par б) Найдите расстояние между этими прямыми, если сторона основания пирамиды равна 2√(2).
  \item Дана правильная шестиугольная пирамида SABCDEF с вершиной S.Точка M−середина бокового ребра CS.
  \par а) Постройте точку пересечения прямой BM с плоскостью ESF.
  \par б) Найдите расстояние между прямыми BM и EF, если сторона основания пирамиды равна 2√(6),а высота пирамиды равна 3√(2).
  \item Основание пирамиды ABCD – равносторонний треугольник ABC,боковое ребро AD перпендикулярно плоскости основания.AD:BC=1:√(2).Точки M и N – середины ребер BC и AB соответственно.
  \par а) Докажите, что угол между прямыми AM и DN равен 60°.
  \par б) Найдите расстояние между этими прямыми, если AB=6√(2).
  \item Основание прямой треугольной призмы ABCA\_1B\_1C\_1−прямоугольный треугольник ABC спрямым углом при вершине A,а боковая грань AA\_1C\_1C−квадрат.
  \par а) Докажите, что прямые CB\_1 и AC\_1 перпендикулярны.
  \par б) Найдите расстояние между этими прямыми, если AC=2;AB\_1=2√(3).
  \item На ребре AA\_1 прямоугольного параллелепипеда ABCDA\_1B\_1C\_1D\_1 взята точка E так, что A\_1E=6EA.Точка T – середина ребра B\_1C\_1.Известно, что AB=4√(2),AD=12,AA\_1=14.
  \par а) Докажите, что плоскость ETD\_1 делит ребро BB\_1 в отношении 4:3.
  \par б) Найдите площадь сечения параллелепипеда плоскостью ETD\_1.
  \item Дан параллелепипед ABCDA\_1B\_1C\_1D\_1 с основаниями ABCD и A\_1B\_1C\_1D\_1.Точки M и N – середины ребер AD и CD соответственно. Точка K лежит на ребре BB\_1,причем B\_1K:KB=1:2.
  \par а) Докажите, что плоскость, проходящая через точки M,N и K делит ребро CC\_1 в отношении 2:7, считая от точки C.
  \par б) Найдите площадь сечения параллелепипеда этой плоскостью, если параллелепипед ABCDA\_1B\_1C\_1D\_1- правильная четырехугольная призма, сторона основания ABCD равна 4√(2),а боковое ребро равно 12.
  \item В правильной шестиугольной пирамиде с вершиной S стороны основания ABCDEF равны 6, а боковые ребра равны 12. Точки K и M – середины ребер SF и SE соответственно.
  \par а) Постройте сечение пирамиды плоскостью BKM.
  \par б) Найдите площадьполученного сечения.
  \item Точка M – середина ребра BD правильного тетраэдра ABCD.Плоскость, проходящая через точку M перпендикулярно ребру AD,пересекает это ребро в точке K,а ребро CD в точке N.
  \par а) Докажите, что N – середина ребра CD.
  \par б) Найдите объем тетраэдра ABCD,если объем пирамиды DKMN равен V.
  \item Плоскость $\alpha$ пересекает боковые ребра $SA$, $SB$ и $SC$ треугольной пирамиды $SABC$ (или их продолжения) соответственно в точках $A_1$, $B_1$, $C_1$. Докажите, что
  \par $\dfrac{V_{SA_1B_1C_1}}{V_{SABC}}=\dfrac{SA_1}{SA}\cdot\dfrac{SB_1}{SB}\cdot\dfrac{SC_1}{SC}$.
  \item Точки M и N-середины ребер соответственно CC\_1 и B\_1C\_1 треугольной призмы ABCA\_1B\_1C\_1 с основаниями ABC и A\_1B\_1C\_1.
  \par а) Докажите, что плоскость BA\_1M делит отрезок AN в отношении 4:3,считая от точки A.
  \par б) Найдите,вкаком отношении плоскость BA\_1M делит объем призмы.
  \item Плоскостьαпроходит через вершину D и центры граней AA\_1B\_1B и BB\_1C\_1C параллелепипеда ABCDA\_1B\_1C\_1D\_1.
  \par а) Докажите, чтоэта плоскость делит ребро BB\_1 в отношении 2:1,считая от вершины B.
  \par б) Найдите объемы многогранников, на которые плоскостьαразбивает параллелепипед, если его объем равен V.
  \item Высота конуса равна 6, а радиус основания равен 8.
  \par а) Докажите, чтонаибольшая площадь сечения конуса плоскостью, проходящей через его вершину равна 50.
  \par б) Найдите расстояние от центра основания конуса до этой плоскости.
  \item Проведены две параллельные плоскости по разные стороны от центра шара на расстоянии 7 друг от друга. Эти плоскости дают в сечении два круга, площади которых равны 9πи 16π.
  \par а) Точка H – ортогональная проекция произвольной точки окружности меньшего круга на плоскость большего. Докажите, что точка H делит проходящий через нее диаметр большей окружности в отношении 1:7.
  \par б) Найдите площадь поверхности шара.
  \item Плоскостьαпроходит через диаметр AB сферы. Через точку B проведена плоскость, касательная к сфере. На этой плоскости взята точка K,причем отрезок KB равен радиусу сферы. Луч AK пересекает сферу в точку M.Через точку M проведена плоскостьβ,перпендикулярная прямой AB.Отрезок CD – общая хорда окружностей сечений сферы плоскостямиαиβ.
  \par а) Докажите, что CD – диаметр окружности сечения сферы плоскостьюβ.
  \par б) Вершина конуса совпадает с точкой B, а окружность основания - с окружностью сечения сферы плоскостьюβ.Найдите объем конуса, если радиус сферы равен 5.
  \item Диаметр окружности основания цилиндра равен 26, образующая цилиндра равна 21.Плоскостьαпересекает его основания по хордам длины 24 и 10.
  \par а) Пусть M и N – середины этих хорд.P−точка пересечения прямой MN с осью цилиндра. Докажите, что расстояния от точки P до плоскостей основания цилиндра относятся как 5:12.
  \par б) Найдите тангенс угла между плоскостьюαи плоскостью основания цилиндра.
\end{enumerate}
\end{document}
